% Gauss, Werke X/1 (1917), native116–117, printed109–110.
\SourceLeaf{116}\ArithmeticHead{109}{R}{GAUSS AN MÖBIUS.}
\NumberedItem{5}
\begin{center}\large\textsc{Gauss} an \textsc{Möbius}.\end{center}
\begin{center}Hochgeschätzter Freund!\end{center}
\noindent\makebox[\linewidth][s]{. . . . . . . . . . . . . . . . . . . . . . . . . . .}\par
\noindent\hspace*{\parindent}Für die jetzt übersandte Abhandlung sage ich meinen verbindlichsten
Dank. Ich habe freilich die genauere Lectüre auf freiere Zeit zurücklegen müssen.
Der Gegenstand gehört einem Felde an, dem Sie schon manche schöne Früchte abgewonnen haben.
Ich möchte Ihrem Nachdenken noch einen analogen Gegenstand empfehlen, zu dessen eigner
Bearbeitung mir schwerlich in langer Zeit Musse gegeben sein wird. Ich meine die verschiedene
Gestaltung und Kreuzung der Linien $T=0$ oder $U=0$ in meiner Schrift von 1799[*)] nach Massgabe
der Ordnung der betreffenden Function. Ich meine die Enumeratio aller verschiedener Fälle
für die Configuration der unendlichen Äste. Anderes damit \textls[120]{verwandtes} hat mich
vielfach beschäftigt, und ich wollte erst zu meiner neulich in der Societ. gehaltenen Vorlesung[**)]
die Darstellung der Hauptmomente jener Untersuchung als 3ten Theil bestimmen; aber ich würde
zur Ausarbeitung dieser Darstellung einer viel grösseren Musse bedurft haben, als mir zu
Gebote gestanden hat.

Ihrem freundlichen Andenken empfehle ich mich angelegentlich und ergebenst
\par\medskip\noindent\hfill C. F. \textsc{Gauss}.\par
\noindent\hspace*{\parindent}Göttingen den 13. August 1849.
\par\bigskip
\begin{center}\small BEMERKUNGEN ZUM VORSTEHENDEN BRIEFWECHSEL.\end{center}
\begingroup\small
\noindent\hspace*{\parindent}Das \textsc{Gauss}archiv besitzt im ganzen zwanzig Briefe von
\textsc{Johann Friedrich Pfaff} an \textsc{Gauss}. Die Briefe von \textsc{Gauss} an
\textsc{Pfaff} sind bis auf einen, vom 21. März 1825, nicht aufzufinden gewesen.
Dieser eine ist abgedruckt in der \textit{Sammlung von Briefen gewechselt zwischen
Joh. Friedr. Pfaff und .\,.\,. andern,} herausgegeben von Dr. \textsc{Carl Pfaff}, Leipzig 1853,
S. 277 ff. Die vorstehenden beiden Briefe von \textsc{Pfaff} werfen ein interessantes Streiflicht
auf die Entstehungszeit von \textsc{Gauss}’ Inauguraldissertation.
\par\smallskip\noindent\rule{.25\linewidth}{.4pt}\par\smallskip
\noindent\hspace*{\parindent}[*) \textit{Demonstratio nova etc.,} Werke III, S. 1; siehe besonders
art. 17, S. 10.]\par
\noindent\hspace*{\parindent}[**) \textit{Beiträge zur Theorie der algebraischen Gleichungen,}
vorgelesen in der Sitzung der Königl. Gesellsch. d. Wissenschaften am 16. Juli 1849,
Werke III, S. 71.]\par
\endgroup
\clearpage
\SourceLeaf{117}\ArithmeticHead{110}{L}{ALGEBRA.\quad BRIEFWECHSEL.}
\begingroup\small
\noindent\hspace*{\parindent}Der Brief von \textsc{Gauss} an \textsc{Drobisch}, der hier nach einer
im \textsc{Gauss}archiv befindlichen Abschrift abgedruckt ist (die Urschrift ist in Privatbesitz),
ist durch die Wertung von Interesse, die \textsc{Gauss} darin den verschiedenen Beweismethoden
für den Fundamentalsatz zu Teil werden läßt. Das Werk von \textsc{Moritz Wilhelm Drobisch},
von dem darin die Rede ist, führt den Titel: \textit{Grundzüge der Lehre von den höheren
numerischen Gleichungen nach ihren analytischen und geometrischen Eigenschaften,}
Leipzig 1834, XXX $+$ 341 S. mit 2 Kupfertafeln; es gibt S. 94 ff. den Beweis der Wurzelexistenz
nach der Inauguraldissertation von \textsc{Gauss}. Die Bemerkung, die \textsc{Gauss} in dem Briefe
in bezug auf seinen zweiten und dritten Beweis macht, ist durch eine auf S. 94 des Buches von
\textsc{Drobisch} befindliche Fußnote veranlaßt, in der es heißt: »Desselben berühmten Geometers
\textit{dem. nova altera etc.,} Gottingae 1816 [Werke III, S. 31] und \textit{demonstratio tertia,}
ibid. eod. ao. [Werke III, S. 57] würden sich theils wegen der Weitläufigkeit des Beweises,
theils wegen der dabei gebrauchten Hülfsmittel hier nicht benutzen lassen«.
Den von \textsc{Cauchy} in der \textit{Analyse algébrique,} 1821, S. 329
(Oeuvres, 2. série, t. III, 1897, S. 274) gegebenen Existenzbeweis *) entwickelt
\textsc{Drobisch} a. a. O. S. 88 ff.; bekanntlich muß aber dieser \textsc{Cauchy}sche Beweis
im Sinne des 4. Einwands, den \textsc{Gauss} 1799 gegen den Beweis von \textsc{d’Alembert}
erhoben hat (Werke III, S. 10), vervollständigt werden, vergl. was \textsc{Gauss} in dem Briefe
[4.] über den Beweis von \textsc{Clausen} bemerkt. \textsc{Gauss} selbst hat in einer Vorlesung,
die er von Neujahr bis Ostern 1840 über die \textit{Theorie der imaginären Größen} gehalten hat,
neben seinem ersten und dritten Beweise, den gedachten Beweis von \textsc{Cauchy} vorgetragen.
Eine Ausarbeitung dieser Vorlesung von \textsc{E. Heine} (vergl. Werke VIII, S. 330, 345)
befindet sich im \textsc{Gauss}archiv. In bezug auf die Bemerkungen, die \textsc{Gauss} über das
Werk von \textsc{Fourier} und seine Anzeige macht, vergl. die weiter unten abgedruckten Stellen
aus Briefen von \textsc{Gauss} an \textsc{Schumacher}. Der von \textsc{Gauss} gebrauchte Ausdruck
\textls[120]{aegri somnia} (S. 106, Zeile 17) stammt aus \textsc{Horaz}, \textit{Ars poëtica} 7,
wo es von einem Buche heißt: »cuius velut aegri somnia vanae fingentur species«.

Der Brief an \textsc{Schumacher} vom 20. Juni 1840 fehlt in der von C. A. F. \textsc{Peters}
veranstalteten Ausgabe des Briefwechsels \textsc{Gauss-Schumacher}. Er ist die Antwort auf den
im dritten Bande, S. 379 abgedruckten Brief Nr. 697 \textsc{Schumachers}.
Die Urschrift befindet sich im \textsc{Gauss}archiv.

Der Brief an A. F. \textsc{Möbius} befindet sich mit noch mehreren andern von \textsc{Gauss}
an denselben Empfänger gerichteten im \textsc{Möbius}archiv der K. S. Gesellsch. der Wissensch.
zu Leipzig **); die Abhandlung, für deren Zusendung \textsc{Gauss} dankt, ist wohl die
\textit{Über die Grundformen der Linien dritter Ordnung} (Abhandlungen der mathem.-phys.
Classe der K. S. Gesellschaft der Wissensch. \textbf{1}, Leipzig 1852***), S. 1,
\textsc{Möbius}’ Werke \textbf{II}, S. 89).
\par\hfill\textsc{Schlesinger.}\par
\smallskip\noindent\rule{.25\linewidth}{.4pt}\par\smallskip
\noindent\hspace*{\parindent}*) Man führt diesen Beweis neuerdings auf J. R. \textsc{Argand}
zurück, siehe die anonym erschienene Schrift \textit{Essai sur une manière de représenter
les quantités imaginaires,} Paris 1806, § 31 und Annales de mathématiques pures et appliquées
\textbf{5}, 1815, S. 201, beides neu herausgegeben von J. \textsc{Hoüel}, Paris 1874.\par
\noindent\hspace*{\parindent}**) Vergl. H. \textsc{Liebmann}, \textit{Aus dem Möbiusarchiv,}
Berichte der math.-phys. Klasse der K. S. Gesellsch. d. Wissensch. zu Leipzig \textbf{32}, 1910,
S. 198; der hier im Auszug abgedruckte Brief ist erst nach dem Erscheinen dieses Aufsatzes
aufgefunden worden, auch ist zu bemerken, daß der Brief von \textsc{Gauss} an \textsc{Möbius}
vom 17. Oktober 1843 sich gegenwärtig im \textsc{Gauss}archiv befindet.\par
\noindent\hspace*{\parindent}***) 1852 ist das Jahr der Fertigstellung des ersten Bandes der
Abhandlungen, die \textsc{Möbius}sche Abhandlung selbst ist bereits 1849 ausgegeben worden.\par
\endgroup
\ClosingRule
\clearpage
